% Options for packages loaded elsewhere
\PassOptionsToPackage{unicode}{hyperref}
\PassOptionsToPackage{hyphens}{url}
%
\documentclass[
]{article}
\usepackage{amsmath,amssymb}
\usepackage{iftex}
\ifPDFTeX
  \usepackage[T1]{fontenc}
  \usepackage[utf8]{inputenc}
  \usepackage{textcomp} % provide euro and other symbols
\else % if luatex or xetex
  \usepackage{unicode-math} % this also loads fontspec
  \defaultfontfeatures{Scale=MatchLowercase}
  \defaultfontfeatures[\rmfamily]{Ligatures=TeX,Scale=1}
\fi
\usepackage{lmodern}
\ifPDFTeX\else
  % xetex/luatex font selection
\fi
% Use upquote if available, for straight quotes in verbatim environments
\IfFileExists{upquote.sty}{\usepackage{upquote}}{}
\IfFileExists{microtype.sty}{% use microtype if available
  \usepackage[]{microtype}
  \UseMicrotypeSet[protrusion]{basicmath} % disable protrusion for tt fonts
}{}
\makeatletter
\@ifundefined{KOMAClassName}{% if non-KOMA class
  \IfFileExists{parskip.sty}{%
    \usepackage{parskip}
  }{% else
    \setlength{\parindent}{0pt}
    \setlength{\parskip}{6pt plus 2pt minus 1pt}}
}{% if KOMA class
  \KOMAoptions{parskip=half}}
\makeatother
\usepackage{xcolor}
\setlength{\emergencystretch}{3em} % prevent overfull lines
\providecommand{\tightlist}{%
  \setlength{\itemsep}{0pt}\setlength{\parskip}{0pt}}
\setcounter{secnumdepth}{-\maxdimen} % remove section numbering
\usepackage{bookmark}
\IfFileExists{xurl.sty}{\usepackage{xurl}}{} % add URL line breaks if available
\urlstyle{same}
\hypersetup{
  pdftitle={Set 1},
  hidelinks,
  pdfcreator={LaTeX via pandoc}}

\title{Set 1}
\author{}
\date{}

\begin{document}
\maketitle

\paragraph{\texorpdfstring{\textbf{1.} Aşağıdaki ifadenin en sade halini
bulunuz.}{1. Aşağıdaki ifadenin en sade halini bulunuz.}}\label{aux15faux11fux131daki-ifadenin-en-sade-halini-bulunuz.}

\paragraph{\texorpdfstring{\textbf{2.} Aşağıdaki ifadenin en sade hâlini
bulunuz.}{2. Aşağıdaki ifadenin en sade hâlini bulunuz.}}\label{aux15faux11fux131daki-ifadenin-en-sade-huxe2lini-bulunuz.}

\paragraph{\texorpdfstring{\textbf{3}. Aşağıdaki önermenin bir totoloji
olduğunu
gösterin.}{3. Aşağıdaki önermenin bir totoloji olduğunu gösterin.}}\label{aux15faux11fux131daki-uxf6nermenin-bir-totoloji-olduux11funu-guxf6sterin.}

\paragraph{\texorpdfstring{\textbf{4.} Aşağıdaki teoremi olmayana ergi
yöntemiyle
ispatlayınız.}{4. Aşağıdaki teoremi olmayana ergi yöntemiyle ispatlayınız.}}\label{aux15faux11fux131daki-teoremi-olmayana-ergi-yuxf6ntemiyle-ispatlayux131nux131z.}

\paragraph{\texorpdfstring{\textbf{5.} Aşağıdaki çarpma işleminin
sonucunu
bulunuz.}{5. Aşağıdaki çarpma işleminin sonucunu bulunuz.}}\label{aux15faux11fux131daki-uxe7arpma-iux15fleminin-sonucunu-bulunuz.}

\paragraph{\texorpdfstring{\textbf{6.} Aşağıdaki sayının onluk tabandaki
karşılığını
bulunuz.}{6. Aşağıdaki sayının onluk tabandaki karşılığını bulunuz.}}\label{aux15faux11fux131daki-sayux131nux131n-onluk-tabandaki-karux15fux131lux131ux11fux131nux131-bulunuz.}

\paragraph{\texorpdfstring{\textbf{7.} Aşağıdaki matrise
göre;}{7. Aşağıdaki matrise göre;}}\label{aux15faux11fux131daki-matrise-guxf6re}

\paragraph{\texorpdfstring{a) \texttt{2A+3B} matrisini
bulunuz.}{a) 2A+3B matrisini bulunuz.}}\label{a-2a3b-matrisini-bulunuz.}

\paragraph{\texorpdfstring{b) \texttt{3B-A} matrisini
bulunuz.}{b) 3B-A matrisini bulunuz.}}\label{b-3b-a-matrisini-bulunuz.}

\paragraph{\texorpdfstring{\textbf{8.} {}}{8. }}\label{section}

\paragraph{\texorpdfstring{\textbf{9.} Aşağıdaki soruyu
cevaplayınız.}{9. Aşağıdaki soruyu cevaplayınız.}}\label{aux15faux11fux131daki-soruyu-cevaplayux131nux131z.}

\paragraph{\texorpdfstring{\textbf{10.} Aşağıdaki denklem sisteminin
çözüm kümesini Cramer yöntemiyle
bulunuz.}{10. Aşağıdaki denklem sisteminin çözüm kümesini Cramer yöntemiyle bulunuz.}}\label{aux15faux11fux131daki-denklem-sisteminin-uxe7uxf6zuxfcm-kuxfcmesini-cramer-yuxf6ntemiyle-bulunuz.}

\end{document}
